% Related-work module included by paper/main.tex.

\section{Related Work and Claim Boundary}
\label{sec:related-work}

Patch intersects mature work on first-class state change, centralized update architectures, effect systems, quantitative typing, capabilities and permissions, event sourcing, operation-based replicated data types, provenance, explicit program changes, translation validation, and proof-carrying evidence. The paper therefore does not argue novelty from any one of these ingredients. The candidate contribution is a particular language-design conjunction: modeled post-creation persistent mutation has a semantic \texttt{Change} route, and the same change representation is reused to derive target-, operation-, and magnitude-aware authority.

\paragraph{First-class and reified state change.}
Plaid makes state change first class and models objects as changing abstract states \cite{sunshine2011plaid}; Worlds reifies program state to control the scope of side effects \cite{warth2011worlds}. These systems rule out any broad claim that Patch invents first-class state transition or reified state. Patch instead studies a deliberately restricted mutation representation whose semantic operations are also consumed by contracts, history, and assurance.

\paragraph{Centralized application-state architectures: Elm, Redux, SAM.}
The Elm Architecture organizes interactive programs around a model and an \texttt{update} function that consumes messages to produce the next model.\footnote{The Elm Guide, ``The Elm Architecture,'' \url{https://guide.elm-lang.org/architecture/}.} Redux likewise makes application state read-only from ordinary clients: state changes are requested by dispatching action objects and reducers compute the next state; its tooling explicitly exploits this discipline for logging, replay, undo/redo, and time-travel debugging.\footnote{Redux documentation, ``Three Principles'' and ``Core Concepts,'' \url{https://redux.js.org/understanding/thinking-in-redux/three-principles} and \url{https://redux.js.org/introduction/core-concepts/}.} The practitioner-oriented State-Action-Model pattern similarly centralizes mutation in a structured propose/accept/learn step.\footnote{J.-J. Dubray, ``Why I No Longer Use MVC Frameworks,'' InfoQ, 2016, \url{https://www.infoq.com/articles/no-more-mvc-frameworks/}; see also \url{https://sam.js.org/}.} These are close conceptual precedents for centralized and explicitly described state transition. Patch does not claim that ``all updates go through one route,'' action objects, reducers, or replay are new. Redux action types and payloads are application-defined, so Redux itself does not assign them a closed magnitude-bearing mutation semantics. Patch's narrower axis is a fixed small operation vocabulary (set/add/remove/clear at the paper level) whose normalized effects can be analyzed uniformly for operation and, in the supported numeric fragment, magnitude.

\paragraph{Event sourcing and event-backed state.}
Event-sourced systems make state-changing events authoritative enough to reconstruct application state, provide history, and support replay; empirical and observability work documents both the benefits and the engineering consequences of treating an event log as state infrastructure \cite{overeem2021eventsourcing,lima2021observability}. This is close architectural prior art to Patch's history/replay motivation. Patch differs in treating the semantic change representation as a language-level mutation interface and in deriving operation/magnitude contracts from that representation; it does not claim the operational or distributed-system properties of event-sourced architectures.

\paragraph{Operation-based CRDTs and explicit deltas.}
Operation-based CRDTs propagate update operations rather than whole replica states, while state-based CRDTs propagate and merge state \cite{shapiro2011crdt}. This is direct prior art for representing updates as operations and for treating operation identity as semantically meaningful. The goals are different: CRDT operation semantics are designed around replicated convergence under communication assumptions, whereas Patch's Paper~1 model is sequential and uses semantic mutation operations for local commit, contracts, and evidence. Operation-based versioning for live executable models likewise records edit and interpreter steps as explicit deltas \cite{exelmans2025operation}. Patch therefore does not claim that operation-based state evolution or execution deltas are themselves new.

\paragraph{Effect systems and effect operations.}
Classical polymorphic effect systems conservatively approximate side effects \cite{lucassen1988effects}. Algebraic accounts identify effect annotations with operation symbols \cite{kammar2012algebraic}, while parametric effect monads \cite{katsumata2014parametric} and effect quantales \cite{gordon2017flow} provide semantic foundations for compositional and order-sensitive effects. Semantic operation labels, interprocedural effect summaries, and effect ordering are therefore not novelty claims for Patch. Patch's restricted design ties its operation summary to the representation used for modeled persistent mutation.

\paragraph{Effects, capabilities, and quantitative authority.}
Effects-as-Capabilities connects effect requirements to capabilities \cite{brachthaeuser2020effects}, and System C further reconciles capability- and effect-based reasoning \cite{brachthaeuser2022boxes}. Graded modal types support quantitative reasoning about resources and effects \cite{orchard2019quantitative}, while dependent-effect work permits grades to depend on program values \cite{kura2026dependent}. Patch does not claim to invent effects-as-capabilities, grades, or value-dependent quantitative effects. Its narrower quantity is the magnitude of a semantic state delta on a target/operation pair, such as an increase bounded by a declared amount.

\paragraph{Permissions, ownership, typestate, and mainstream controlled mutation.}
Borrowing permissions provide static reasoning about aliases and access rights \cite{naden2012borrowing}; Mezzo uses affine and duplicable permissions to control ownership, aliasing, and typestate changes \cite{pottier2013mezzo}. Recent work on revocable capabilities makes capability availability flow-sensitive \cite{jia2026revocable}, and InvalML tracks permanent and temporary invalidation with a type-and-effect system \cite{gao2025invalidation}. Rust is also an important mainstream counterexample to any suggestion that controlled mutation routes are exotic: mutable access is expressed through \texttt{\&mut} references and constrained by the borrow rules.\footnote{The Rust Programming Language, ``References and Borrowing,'' \url{https://doc.rust-lang.org/book/ch04-02-references-and-borrowing.html}.} Patch's concern is different from alias control or memory safety: a permitted mutation is further classified by target, semantic operation, and supported magnitude.

\paragraph{Hoare, dependent, and refinement-style verification.}
Hoare Type Theory integrates pre/post specifications with mutable higher-order state \cite{nanevski2008htt}. F* supports effectful verification through dependent types and predicate-transformer semantics \cite{swamy2016fstar}, and Dijkstra monads derive weakest-precondition reasoning for user-defined effects \cite{ahman2017dijkstra}. Such systems can express state relations at least as precise as a bounded Patch update. Patch therefore does not claim unique expressibility for a rule such as ``the new score is no more than ten above the old score.'' The design question is whether a smaller mutation-centered vocabulary can infer useful summaries from the update representation itself.

\paragraph{Explicit changes, edits, and patch theory.}
COPE advocates change-oriented programming environments \cite{dig2016cope}; Edit Transactions group live program edits into dynamically scoped change sets \cite{mattis2017edittransactions}. Higher-order change calculi support incremental computation \cite{cai2014changes}, while edit lenses \cite{hofmann2012edit} and homotopical patch theory \cite{anguilhomotopical2016} provide rich theories of edits and repository patches. These works prevent any claim that programming with explicit changes or formalized patches is new; their changed entities and primary goals differ from Patch's runtime application-state commit route.

\paragraph{Provenance and causal explanation.}
Database provenance distinguishes why-, how-, and where-provenance and supports applications including debugging, annotation propagation, and explanation \cite{cheney2009provenance}. Whyline-style debugging likewise provides causal queries over program behavior \cite{ko2008whyline}. Patch uses the term \emph{provenance} only for information made available by its explicit mutation history. It does not claim the richer provenance semantics of database provenance or the causal power of Whyline.

\paragraph{Assurance by checking unverified producers.}
Translation validation \cite{necula2000translation} and Proof-Carrying Code \cite{necula1997pcc} are direct prior art for checking evidence produced by components that are not themselves verified. Patch's generated Lean certificates and independent validators reuse this assurance pattern. The runtime bridge is therefore supporting assurance, not a claim that checking unverified evidence producers is novel.

\paragraph{Claim boundary.}
The related work is intentionally discussed in prose rather than in a yes/no feature matrix. Several cited systems can encode properties that Patch exposes directly, and a matrix whose columns are derived from Patch's own architecture would privilege Patch by construction. The defensible claim is narrower: Patch makes a semantic \texttt{Change} object the modeled post-creation persistent-mutation interface and reuses that object to derive a restricted target/operation/magnitude contract vocabulary. Whether that conjunction is useful is evaluated by the controlled cases and the bounded public-code stress test; the paper does not assert firstness for centralized mutation, action/event logs, explicit deltas, quantitative effects, capabilities, provenance, or controlled mutation.

\paragraph{Future-work boundary.}
Relational multi-target ChangeSets, least-authority capability inference, certified foreign-state adapters, temporal authority, causal debugging, and safe parallelism remain separate follow-on research questions. The mutation-shape stress test can motivate those directions without turning them into Paper~1 contributions.
